\section{The World Economy and National Accounts}


\begin{frame}{The World is large---and all we have}
    \mycomment{The World Economy is the only `closed economy' we have. 
    
    There is production and consumption, and there is global growth---
    the world economy is not a zero sum game. 
    
    There are also fundamental balances: every buyer has a seller, every import is also an export.
    }
    \showmybitmap[https://earth.google.com/]{../shared/bitmaps/world-googleearth.png}
\end{frame}

\begin{frame}{Orders of Magnitude}
    \mycomment{
    Since we are dealing with large numbers (of monetary terms), we need to have some prices and figures to put them into perspective.
    
    How much is an oil tanker worth? How much its load? What is the GDP of your home country? How much is the grain worth traded per day?
    }
    \showmybitmap{../shared/bitmaps/joyofmoney.jpg}
\end{frame}
\begin{frame}{Orders of Magnitude}
    \mytable{
       \begin{tabular}{lrrll|l}
        Sc. & Acc. & IT & engl. & other (D) & \textit{example} \\
        \toprule
         $10^3$ & TEUR & K & thousand &  & \phantom{laptop/computer}\\
        $10^4$ & & 10K & &  & \\
        $10^5$ & & 100K & & \emph{lakh} & \\
        \midrule
        $10^6$ & MEUR& 1M & million & \emph{Millionen} &  \\
        $10^7$ & & 10M & &  \emph{crore} & \\
        $10^8$ & & 100M & &  & \\
        \midrule
        $10^9$ &  & 1G & billion & \emph{Milliarden}  & \\
        $10^{10}$ & & 10G & &  & \\
        $10^{11}$ &  & 100G & &  & \\
        \midrule
        $10^{12}$ &  & 1T & trillion & \emph{Billionen} & \\
        \bottomrule
    \end{tabular}
    }
\end{frame}

\begin{frame}{Orders of Magnitude}
    \mycomment{
    To have an idea about the meaning of (large, economic) numbers, it makes sense to memorise a couple of them. For instance, the GDP of your home country (and maybe of one or two larger economies) is a good standard candle. The budget of your city council could be another.
    
    Every now and then, it may make sense to compare big numbers to unrelated numbers. For example, there are (more than) 31 million seconds in a year, and there are 100 billion stars in the galaxy. Just how many cars can you buy for the price of one oil tanker? How many nursing homes can you build for the price of one moon mission?
    }
    \mytext{
    How much...
    \mylonglist{
    \begin{itemize}
        \item costs a car, a house?
        \item the whole super market?
        \item prime estate in London?
        \item the World Trade Center?
        \item a Panamax oil tanker, and its content?
        \item turnover in Bremerhaven per day?
        \item A container shipping Shanghai-Hamburg?
        \item The Artemis II moon mission?
        \item The companies Siemens, Saudi Aramco?
        \item US government interest payment?
        \item US government defense budget?
    \end{itemize}
    }
    }
\end{frame}

\begin{frame}{Stock Quantities vs. Flow Quantities}
    \mycomment{
    Numbers need to be put in relation (to other numbers), so that we can make sense of them. Small percentages of large numbers are large numbers also. If you double a number of a certain order of magnitude, you get a new number in the same order of magnitude. 

    Most importantly: are we dealing with absolute quantities, or with the growth of quantities?
    }
    \mytext{
    \begin{description}
        \item[Stock quantities] describe how much there is of something at a present time. 
        \item[Flow quantities] describe how much is added or subtracted per time frame (usually a year).
        \item[Orders of Magnitude] are often more important than precise numbers,
        \item[Growth Rates] are connecting stock and flow quantities; discern 1st and 2nd order derivatives to describe trends.
    \end{description}
    }
\end{frame}

\begin{frame}{Flow vs. Derivatives / Marginal Rates}
    \mycomment{
        To make numbers speak---for us, to understand numbers properly---we must always put them into perspective. So for us, growth rates are typically more interesting than absolutes; and absolutes are intersting when put into relation to other absolutes. 
    }
    \showmybitmap{../shared/bitmaps/only-a-sith-deals-in-absolutes.png}
\end{frame}

\begin{frame}{Flow vs. Derivatives / Marginal Rates}
    \mytext{
     \begin{itemize}
         \item Flow rates: how much more \textit{per period}
         \item Marginal rates: how much more \textit{per unit}
     \end{itemize}
     Mathematically, \textit{marginal rates} are the derivative of the function:
     \begin{equation}
         f'(x) = \frac{df}{dx}
     \end{equation}
     Since marginal rates are important for optimisation, they are more relevant for decision making than stock or flow values.
    }
\end{frame}



\begin{frame}{Discounting}
    \mytext{
    The growth rate of anything is often more relevant than the sum. 
    \begin{equation}
        K_t = K_0 \left(1+\frac{p}{100} \right)^t
    \end{equation}

    Investors will put their money where they receive highest returns, and shy away from lower or negative returns. 
  

    }
    
\end{frame}

\begin{frame}{Rule of 72}
    \mytext{

  
    % For practical reasons, you can use the `rule of 72'.
    With this rule of thumb you can estimate how many years it takes for a given start capital to double, given a specific growth rate:
    \begin{equation}
        p \approx \frac{72}{\Delta t} \quad \Leftrightarrow \quad \Delta t=\frac{72}{p}
    \end{equation}

    Stock at a growth rate of 5\% will double in approximately in $\Delta t\left(p=5\right) \approx \frac{72}{5}=13$ years. 
    }
    
\end{frame}


\begin{frame}{Elasticity of Demand}
    \showmytikz[1]{../shared/tikz/tikz-elasticity.tex}

    \mycomment{
    A demand curve can be more or less steep. A flat demand curve means that demand (quantity) reacts very sensitive to price change; a steep demand curve means that demand (quantity) is little sensitive to price change. 


    This is called `price elasticity of demand':

    \begin{equation}
        \eta = \frac{
            \frac{d Q}{Q}
        }{
            \frac{d P}{P}
        }
    \end{equation}

    }    
\end{frame}

\begin{frame}{Elasticity of Demand}
    \showmytikz[1,2]{../shared/tikz/tikz-elasticity.tex}

    \mycomment{


    (price-)\textbf{inelastic demand:} e.g. energy 

    (price-)\textbf{elastic demand:} e.g. carousel tickets 

    Elasticity relates to the availability of substitutes.
    }    
\end{frame}

\begin{frame}{Elasticity of Demand}
    \showmytikz[1]{../shared/tikz/tikz-elasticitytwo.tex}
    \mycomment{
        So in general, a flat curve is inelastic, and a steep curve elastic. In reality, demand curves are not linear.
    
        Note that a linear demand curve does not have constant elasticity. That is because the relation the absolute values of $P,Q$ is different at different parts of the slope.

        \begin{equation}
            \eta = \frac{
                \frac{\delta Q}{Q}
            }{
                \frac{\delta P}{P}
            }
        \end{equation}

        Between $A,B$ a very small percentage change of the price $P$ has resulted in a larger percentage change of the quantity $Q$.

        At high elasticity, a tiny price change will have a large quantity effect. 

        And even a large quantity change will change the price just a little. 
        
    }
\end{frame}

\begin{frame}{Elasticity of Demand}
    \showmytikz[1,2]{../shared/tikz/tikz-elasticitytwo.tex}
    \mycomment{
         Between $C,D$ a larger percentage change of the price $P$ has resulted in a small percentage change of the quantity $Q$.

        At low elasticity, small quantity changes affect large price changes, while large price changes affect quantity little.
        
    }
\end{frame}

\begin{frame}{Elasticity of Demand}
    \showmytikz[1,2,3]{../shared/tikz/tikz-elasticitytwo.tex}
    \mycomment{
         At $U$ we have unity elasticity: a 1\% change in price will result in a 1\% change in quantity.

         It follows that maximum revenue is at the point of unity elasticity (ceteris paribus).

         Elasticity depends on the slope of the demand curve, but also on the location in the demand curve!

         Note that demand curves are in reality not linear. They may be more or less linear at the point of equilibrium, so we are typically concerned with the elasticity \textit{at that point.}
    }
\end{frame}

\begin{frame}{Elasticity of Demand}
    \showmytikz[1,2]{../shared/tikz/tikz-elasticity-oil.tex}
    \mycomment{
         The Iran war 2026 resulted in closure of the straight of Hormus, through which approximately $20\%$ of crude oil is transported. Oil price almost doubled. So here, $\eta=\frac{0.2}{0.5}=0.4$.

         Energy is the most quoted example of a commodity with very low elasticity of demand. How much less do you guess do people drive / use oil because of a $50\%$ increase in price? Well, apparently only about $20\%$ less. 

         At the time of writing (April 12th 2026) the longer term effects have not yet been visible. It is insofar not sound  to set $\Delta Q=20\%$.

         Also note that short-term elasticity and long-term elasticity differ, as buyers may be able to find substitutes after a while.
         
        Note that the mathematical definition of elasticity depends on marginal changes ($\delta x$ or $dx$). Here, we have taken absolute changes ($\Delta x$). Sometimes, one is e.g. denoted by $\eta$ and one by $\epsilon$.

        How have prices of diesel and petrol reacted relative to crude?
    }
\end{frame}
\begin{frame}{Elasticity of Demand}
    \mycomment{
    Cross-price elasticity measures the quantity change of a product when the price of \textit{another} product changes. It is an important tool to analyse complements and substitutes, precursor products, and by-products.

    For example, if plastic is a by-product of crude oil, the price of plastic will react when crude prices change. Silver is a by-product of copper mining. The price of raw materials will change depending on industrial demand, so depend on where and how they are used.

    }
    \mytext{
    Price elasticity of demand:
    \begin{equation}
        ED=\frac{
            \frac{dx}{x}
        }{
            \frac{dp_x}{p_x}
        }
    \end{equation}
    Cross-price elasticity of demand:
    \begin{equation}
        XED=\frac{
            \frac{dx}{x}
        }{
            \frac{dp_y}{p_y}
        }
    \end{equation}
    \begin{equation}
        \begin{array}{rl}
            XED_{x,y} <0  & \quad \rightarrow x,y\quad \text{are substitutes}\\
            XED_{x,y} >0 & \quad \rightarrow x,y\quad \text{are complements}\\
        \end{array}
    \end{equation}
    
    }
\end{frame}


\begin{frame}{Macroeconomics: Gross Domestic Product}
    \mycomment{
    Eq.~\ref{eq:fund} is called the fundamental equation of Macroeconomics representing the expenditure approach. These identities, most importantly $S=I$ in a closed economy, are fundamental, as they are account-based.
    }
    \mytext{
    Sum of everything produced / used in a country
    \begin{equation}\label{eq:fund}
        Y = C + I + Ex - Im + G
    \end{equation}

    Everything produced ($Y$) is either
    \begin{itemize}
        \item consumed ($C$) or 
        \item invested ($I$, equals savings $I=S+ \Delta K$), or
        \item exported ($Ex$) or
        \item not produced, but imported ($Im$) or
        \item state expenditure ($G$).
    \end{itemize}

    Note that GDP is a flow quantity.
    }

\end{frame}

% \section{National Accounts and Balance of Payments}

\begin{frame}{National Accounts}
    \mytext{
    The System of National Accounts (SNA) is a system of book-keeping on the national level. \cite{SysNatAcc2025} \cite{lequillerblades2006}

    There are \textbf{current accounts} and \textbf{accumulation accounts.} Both of them are comparable to income statements, not to balance sheets.
    }
    \mycomment{Historically, there are variations and also slight differences on the national level. There has been a wide degree of standardisation, particularly after WWII. This comes from institutions like the IMF or the EU, which require standardised reporting.}
\end{frame}

\begin{frame}{National Accounts}
    \mycomment{Current accounts express flow values (within a particular time frame, typically a year). One example is GDP (per year).

    One important current account is the \textbf{trade balance: the account of imports of goods and services.}
    }
    \mytext{

     \textbf{Current accounts}
        \begin{itemize}
            \item production accounts
            \item primary distribution of income
            \item secondary distribution of income
            \item \textbf{imports and exports of goods and services}
            \item use of income accounts
        \end{itemize}

    }
\end{frame}
\begin{frame}{National Accounts}
    \mycomment{Accumulation accounts express the change of assets held by institutional units. The most important of this is the \textbf{capital account}. The part of the (overall) capital account that crosses borders appears as sub-account in the  Balance of Payments. 

    % On top of that, and complementary of that, \textbf{additional accounts} are held. These are sometimes the result of double book keeping and highlight other aspects of the economy.

    % \textbf{Balance sheets} per institutional unit, sector, group
    
    % \textbf{Goods and services accounts}
    
    % \textbf{Accounts for the rest of the world}


    
    }
    \mytext{
    \textbf{Accumulation accounts}
        \begin{itemize}
            \item \textbf{capital account}
            \item \textbf{financial account}
            \item other changes in assets account
            \item revaluation account
        \end{itemize}
    
 

    %% https://ec.europa.eu/eurostat/statistics-explained/index.php?title=Building_the_System_of_National_Accounts_-_basic_concepts
    }
\end{frame}






\begin{frame}{Balance of Payments, Simplified}
    \showmytikz{../shared/tikz/tikz-boptikz}
    \mycomment{
    The Current Account appears on the debit side, the Financial and Capital Accounts appear on the credit side of the Balance of Payments.

    Depending on whether the saldo of the Trade Balance is positive or negative, the Trade Balance will appear on the left or right of the BoP. The same applies for the other Accounts.

    \textbf{The Balance of Payments is always level.} There can be deficits in the Current Account, or in the Financial and Capital Accounts. The Reserve Account serves as saldo. It shows how much (foreign) currency the central bank holds.

    Contrary to what the names suggest, the Balance of Payments is not a balance, but shows \textbf{flow quantities: inflows, outflows over a period.}
    }
    
\end{frame}



\begin{frame}{Imports and Exports in the BoP}
    \showmytikz{../shared/tikz/tikz-bop-entries}
    \mycomment{
    Export of goods and services is either payed for in cash ($\Delta R$), or made possible through an export of a liability. Thus normally, exports of goods correspond to export of capital.

    A deficit in the current account is thus always mirrored by a deficit in the capital account, as the export of goods is accompanied by the export of a liability, which is an import of capital.

    Foreign capital inflow makes current account deficits possible. 

    Note that currency reserves also include foreign government securities, in particular those of the US (USTs).
    }

\end{frame}
\begin{frame}{Balance of Payments in Balance}

    \mycomment{
        There is sometimes confusion: a trade deficit or surplus means an imbalance in the current account, or the trade balance as part of the current account.

        It follows that imports not covered by exports ($Im>Ex$) can only appear when capital import happens ($K_{Im}>K_{Ex}$; that is the credit the exporter gives the importer.

        The Reserve Account consists of cash, but also of e.g. government securities. This is often UST (T-Bills), in particular in USD.

        The rest of the world is financing the US trade deficit through capital import, often by buying highly liquid USTs. Weaker currencies, states or government do not hold that privilege.
    }
        
    \mytext{
        The BoP is \textbf{always, by definition, in balance.} 
        
        \begin{equation}
            Ex-Im-K_{Ex}+K_{Im}-\Delta R = 0
        \end{equation}

        To see this, look at the sub-balances:
        
        \begin{equation}
            \begin{array}{lllll}
             Ex - Im & - & (K_{Ex} - K_{Im}) & = & \Delta R  \\
             \text{\scriptsize (Current A.)} & - & \text{\scriptsize (Fin \& Cap. A.)} & = & \text{\scriptsize (Reserve A.)  }
            \end{array}
        \end{equation}


    }
\end{frame}

\begin{frame}{Balance of Payments (IMF)}
    \mycomment{
    This tree represents the Balance of Payments (account) as defined by the IMF \cite{imfbpm7}. The BoP is the account for the Rest of the World, that is $IM, EX$ in the fundamental equation of Macroeconomics.
    }
    \showmytikz{../shared/tikz/tikz-bop.tex}
\end{frame}


\begin{frame}{GDP in USD, 2025}
    \mycomment{
    GDP of countries are the main measure for the economic importance of a country. However, it does not count stocks. There could be countries (Swizerland? UK? Saudi Arabia?) that hold many assets, but produce relatively little. The income generated (for example in the US) is counted where it is produced - independent of the fact that it may end up being owned abroad. 
    }

    
    \mytable{
        \begin{tabular}{rlrl}
            \toprule
            \def\myblock{fill=black}
            & & Trillion USD &  \\
            & World & 110.0 &  \\
           \midrule
            1st & USA  & 30.6   & \myblock[0.4ex]{30.6}\\
                & EU   & 21.1   & \myblock[0.4ex]{21.1}\\
            2nd & China & 19.4  & \myblock[0.4ex]{19.4}\\
            3rd & Germany & 5.0 & \myblock[0.4ex]{5.0}\\
            4th & Japan & 4.3   & \myblock[0.4ex]{4.3}\\
            5th & India & 4.1   & \myblock[0.4ex]{4.1}\\
            $\vdots$ & $\vdots$ & $\vdots$ & \\
            \bottomrule
        
        \end{tabular}

        Note: 1 Trillion $= 10^{12}$
    
        Source: \cite{IMFdatamapper}
    
    }
\end{frame}

\begin{frame}{GDP Growth Rates, 2026}
    \mycomment{
    Growth rates are interesting, because they give a hint how the absolute ranking may change. Refer to the `rule of 72'. At a growth rate of 4\%, any value will approximately double after $72/4=36$ periods.
    }
    \mytable{
    \begin{tabular}{rlrl}
        \toprule
        & & Trillion USD     &  \\
        & World & 110.0      &  \\
       \midrule
        1st & USA  & 2.1     & \myblock[2ex]{2.1}\\
            & EU   & 1.4     & \myblock[2ex]{1.4}\\
        2nd & China & 4.2    & \myblock[2ex]{4.2}\\
        3rd & Germany & 0.9  & \myblock[2ex]{0.9}\\
        4th & Japan & 0.6    & \myblock[2ex]{0.6}\\
        5th & India & 6.2    & \myblock[2ex]{6.2}\\
        $\vdots$ & $\vdots$ & $\vdots$ &  \\
        \bottomrule
    \end{tabular}

    Source: \cite{IMFdatamapper}
    
    }
\end{frame}

\begin{frame}{GDP in USD per Capita 2025}
    \mycomment{The absolute GDP does not tell us about the standard of living or grade of development of a country. There are countless measures, but the easiest to handle is GDP per capita.
    
    This does of course not take into account how evenly or unevenly this total income may be distributed. There may be much inequality, and much poverty, in the `richest' of nations.
    }
    \mytable{
        \begin{tabular}{lrl}
            \toprule
            Country/Region & TUSD/a & \\
            \midrule
            Switzerland & 110       & \myblock[0.2ex]{110}\\
            USA &  89               & \myblock[0.2ex]{89}\\
            Germany & 59            & \myblock[0.2ex]{59}\\
            Europe & 42             & \myblock[0.2ex]{42}\\
            Japan & 34              & \myblock[0.2ex]{34}\\
            Poland & 28             & \myblock[0.2ex]{28}\\
            Mexico & 14             & \myblock[0.2ex]{14}\\
            \textbf{World} & 15     & \myblock[0.2ex]{15}\\
            China & 14              & \myblock[0.2ex]{14}\\
            India & 3               & \myblock[0.2ex]{3}\\
            \bottomrule
        \end{tabular}
        
         Source: \cite{IMFdatamapper}
    }
\end{frame}

\begin{frame}{Trade as Part of Global GDP}
    \mycomment{Trade is a large part of global GDP. That is, a large part of what is produced world wide is traded. 
    
    That does of course not measure precisely the contribution of trade to GDP (and GDP growth), which is hard to measure.
    }
    \mytable{
    \begin{itemize}
        \item Global trade 2024: ca. 33 Trillion USD
        \item 8-9 trillion USD services
        \item 24 trillion USD goods
    \end{itemize}
    Source: \href{https://unctadstat.unctad.org/datacentre/dataviewer/US.TradeMerchTotal}{UNCTAD}

    As global GDP stands at 110 Trillion USD, trade is about \textbf{30\%.}
    }
\end{frame}

\begin{frame}{Trade Indicators}
\mycomment{
    The (trade) openness of a country is measured by setting the sum of exports and imports relative to gross domestic product.
    
    It does not, however, tell us what type of goods are imported and exported: raw materials, or consumer goods, or energy? And are these all imported or exported to a plethora of, or are there dependencies on singular trading partners?
    
    One can of course also put just imports or just exports in relation to GDP. This gives you the Import or Export Penetration Rate.

    The Dependency Rate measures how much of a good $X$ (raw material, energy) is imported relative to how much of it is consumed locally. 

    }
\mytext{

    Trade Openness:
    \begin{equation}
        TO = \frac{\left|{EX}\right|+\left|{IM}\right|}{{Y}}
    \end{equation}  

    Import Penetration Rate:
    \begin{equation}
        IPR = \frac{IM}{Y}
    \end{equation}

    Dependency Rate:
    \begin{equation}
        DR_X=\frac{IM_X}{C_X}
    \end{equation}
}
\end{frame}
\begin{frame}{Trade Indicators}
\mycomment{
Openness varies widely from country to country. But larger countries are of course `less open' than small countries. This does not say anything about their trade policy, but only about the part of GDP that is traded. 

Remember that trade openness counts imports and exports, thus every cross-border trade is counted twice. To measure the influence on GDP, one has to look at the trade balance in accordance to the fundamental equation of macroeconomics (eq.~\ref{eq:fund}).

}
\mytable{
    \begin{tabular}{lrl}
        \toprule
         Brasil & 38.8               & \myblock[0.18ex]{38.8}\\
         Chile & 75.2                & \myblock[0.18ex]{75.2}\\
         China & 37.4                & \myblock[0.18ex]{37.4}\\
         Congo Rep. & 98.2           & \myblock[0.18ex]{98.2}\\
         EU & 103.7                  & \myblock[0.18ex]{103.7}\\
         Germany & 88.8              & \myblock[0.18ex]{88.8}\\
         High Income & 67.4          & \myblock[0.18ex]{67.4}\\
         Hong Kong & \textbf{352.9}          & \\
         Indonesia & 45.5            & \myblock[0.18ex]{45.5}\\
         Least developed (UN) & 51.2 & \myblock[0.18ex]{51.2}\\
         Mexico & 88.8               & \myblock[0.18ex]{88.8}\\
         UK & 69.2                   & \myblock[0.18ex]{69.2}\\
         USA & 27.3                  & \myblock[0.18ex]{27.3}\\
         World & 62.4                & \myblock[0.18ex]{62.4}\\
         \bottomrule
        \bottomrule
    \end{tabular}
    
    Data for 2022 \cite{worldbankdatabank}
}
\end{frame}

\begin{frame}{Imports Ratio to GDP}
    \mycomment{
    As presented by Kristalina Georgieva, Managing Director, IMF on April 17th, 2025 \cite{IMFgeorgieva2025}. Large economies are typically much less open than smaller ones. 
    
    Protectionism affects smaller countries more than larger ones.
    }
    \showmybitmap[https://www.imf.org/en/news/articles/2025/04/17/sp041725-spring-meetings-2025-curtain-raiser]{../shared/bitmaps/6-curtain-raiser-charts-sms-final.jpg}
\end{frame}


